\section{iGPT: ¿puede el preentrenamiento autoregressive transferirse a Pixels?}

Si el núcleo del preentrenamiento generative es modelar la distribución de una secuencia, una imagen también puede cuantizarse, aplanarse y predecirse en orden de píxeles o de code. El problema es que la imagen tiene estructura local bidimensional y el raster order es una elección arbitraria; una causal factorization razonable para el lenguaje no es necesariamente el sesgo inductivo más eficaz para la visión.

\subsection{Convertir la imagen en una secuencia}

Esta sección convierte primero la imagen en una secuencia que un decoder pueda procesar. iGPT representa la imagen como tokens de píxel y se entrena con next-pixel prediction; el modelo puede generar completion y también usar su hidden representation para classification probing.

\lecturefigure{slide-22-image-gpt-question.jpg}{Image GPT pone a prueba si la misma receta autoregressive se extiende a las imágenes.}{Video oficial, 23:42; Chen et al. 2020.}

\begin{knowledgebox}{Lectura de la figura: a la izquierda, la pregunta; a la derecha, la factorization}
La imagen se despliega en un orden fijo y el modelo predice el siguiente píxel o valor discreto. Así puede reutilizarse directamente un Transformer decoder-only, pero la adyacencia bidimensional se convierte en un orden unidimensional. El éxito muestra que la architecture es general; la brecha en eficiencia y en sample quality muestra que el bias del dominio sigue siendo importante.
\end{knowledgebox}

\subsection{Completions: la versión visual de la distribución condicional}

A continuación se usa la completion para observar la distribución condicional. Dada la mitad superior de una imagen, el modelo genera la mitad inferior; la diversidad de salidas indica que existen varias plausible continuation, no que el modelo «recupere la única imagen verdadera». Antes de leer la figura conviene preguntar qué conserva la restricción de entrada, qué atributos siguen siendo inciertos y si las distintas muestras cubren realmente varios modos razonables.

\lecturefigure{slide-23-igpt-completions.jpg}{Las completion de iGPT muestran la incertidumbre multimodal de la generación visual condicional.}{Video oficial, 25:31.}

\begin{knowledgebox}{Lectura de la figura: la respuesta correcta no son los píxeles de la imagen original}
Si la mitad superior de un rostro puede corresponder a varias mitades inferiores, la ground truth píxel por píxel es solo una muestra. La evaluación debe considerar la perceptual plausibility, la global consistency y la diversity. Usar solo el pixel MSE penaliza generaciones razonables pero distintas, y también puede favorecer promedios borrosos.
\end{knowledgebox}

\subsection{Feature learning}

Después se examina si el objective generativo aprende feature útiles para clasificación; este es el punto central del transfer loop. La variación de la probe accuracy según la capa o el tamaño del modelo puede mostrar dónde es más separable la representation.

\lecturefigure{slide-24-igpt-feature-learning.jpg}{iGPT evalúa a la vez las generative samples y las learned visual features.}{Video oficial, 26:24.}

\begin{knowledgebox}{Lectura de la figura: la calidad de generación y la calidad de representación son dos ejes}
La tabla compara modelos, capas o evaluation protocol. Un modelo puede generar imágenes vistosas sin ofrecer el mejor linear probe, o al revés. El representation learning requiere congelar las feature y unificar el classifier y la data split, para no confundir la probe capacity con la backbone quality.
\end{knowledgebox}

\subsection{Resumen de la sección}

iGPT demuestra que el preentrenamiento autoregressive puede transferirse a pixels y producir a la vez generation y representation. También pone de manifiesto la importancia del sequence order, del sesgo visual y del protocolo de evaluación.
